\upshape
\emph{Idea.} The sign of $\Delta_\psi(e)-b$ is that of $\bar\psi(e)-\psi$, and $\bar\psi$ is increasing; since every equilibrium silent price lies between its values at $e(1)$ and $\bar\theta$, comparing $\psi$ with $\bar\psi$ at the two ends classifies every equilibrium.

\emph{Step 1 (two facts when $e^{\rm def}>1$).} Let $e^{\rm def}>1$. By the identity in the proof of Proposition~\ref{prop:att_seller_types},
\begin{equation*}
    \Delta_\psi(e)-b\;=\;b\,\big(e^{\rm def}-e\big)\big(\bar\psi(e)-\psi\big),\qquad
    \bar\psi'(e)\;=\;\frac{e^{\rm def}-1}{(e^{\rm def}-e)^2}\;>\;0,
\end{equation*}
so $\bar\psi(e(1))<\bar\psi(\bar\theta)$. Any equilibrium with $\sigma_1^*=0$ has $\Delta^*\in[\Delta_\psi(e(1)),\Delta_\psi(\bar\theta)]$ by~\eqref{eq:att_premium_trunc}.

\emph{Step 2 (the three cases when $e^{\rm def}>1$).}
\begin{itemize}
    \item If $\psi<\bar\psi(e(1))$, then $\Delta^*\ge\Delta_\psi(e(1))>b$. Any equilibrium with $\sigma_1^*>0$ has $\Delta^*\ge b+c_0+c_1>b$.
    \item If $\psi\ge\bar\psi(\bar\theta)$, then $\Delta^*\le\Delta_\psi(\bar\theta)\le b$. Moreover $\bar g_1(\sigma_1)\le\bar g_1(1)=\Delta_\psi(\bar\theta)-b\le0<c_0+c_1$ by Lemma~\ref{lem:etilde_monotone} and $e(1,1)=\bar\theta$, so no equilibrium has $\sigma_1^*>0$.
    \item Otherwise $\Delta_\psi(e(1))\le b<\Delta_\psi(\bar\theta)$, so $b$ lies inside the band. By~\eqref{eq:att_premium_trunc}, $\Delta^*>b$ if and only if $c_0>b$.
\end{itemize}

\emph{Step 3 (the case $e^{\rm def}\le1$).} Then $\tilde e\le1$ at every pool, and the second case applies by the same argument.

\emph{Step 4 (the model of Section~\ref{sec:equilibrium}).} With $\psi=1$, $\bar\psi(\bar\theta)<1=\psi$. And $0<e(1)\le e^*\le\bar\theta<1$ gives $0<\Delta^*<b$.
